% ============================================================
%  IEEE Transactions — Kill-Chain MaxSAT Honeypot Placement
%  Student-Level Writing + Full Formulation Examples — v4
% ============================================================
\documentclass[journal]{IEEEtran}

\usepackage{amsmath,amssymb,amsfonts}
\usepackage{algorithmic}
\usepackage{algorithm}
\usepackage{array}
\usepackage{graphicx}
\usepackage{textcomp}
\usepackage{xcolor}
\usepackage{booktabs}
\usepackage{multirow}
\usepackage{url}
\usepackage{hyperref}
\usepackage{cite}
\usepackage{balance}
\usepackage{pifont}
\usepackage{mdframed}
\usepackage{subcaption}

\hypersetup{colorlinks=true,linkcolor=blue,citecolor=blue,urlcolor=blue}

% ── check/cross marks ────────────────────────────────────────
\newcommand{\cmark}{\ding{51}}
\newcommand{\xmark}{\ding{55}}

% ── coloured example boxes ───────────────────────────────────
\definecolor{exbg}{RGB}{240,248,255}
\definecolor{exborder}{RGB}{0,102,204}
\newmdenv[
  backgroundcolor=exbg,
  linecolor=exborder,
  linewidth=1.5pt,
  leftmargin=2pt, rightmargin=2pt,
  innerleftmargin=5pt, innerrightmargin=5pt,
  innertopmargin=4pt, innerbottommargin=4pt
]{examplebox}

% ── macros ───────────────────────────────────────────────────
\newcommand{\ie}{\emph{i.e.}}
\newcommand{\eg}{\emph{e.g.}}
\newcommand{\etal}{\emph{et~al.}}
\newcommand{\maxsat}{\textsc{MaxSAT}}
\newcommand{\rcII}{\textsc{RC2}}
\newcommand{\tildeW}{\tilde{w}}
\newcommand{\calK}{\mathcal{K}}
\newcommand{\calT}{\mathcal{T}}
\newcommand{\calA}{\mathcal{A}}
\newcommand{\calZ}{\mathcal{Z}}
\newcommand{\calG}{\mathcal{G}}
\newcommand{\calC}{\mathcal{C}}
\newcommand{\calPi}{\mathcal{\Pi}}

% ============================================================
\begin{document}

\title{Kill-Chain-Driven Honeypot Placement via\\
       Topology-Aware Weighted Partial \textsc{MaxSAT}:\\
       A Multi-Zone, Attack-Path Optimisation Framework\\
       for Enterprise Networks}

\author{
  \IEEEauthorblockN{Tu Nguyen (PhD Student) and
                    Professor Haadi Jafarian}\\
  \IEEEauthorblockA{Department of Computer Science,
    University of Colorado Denver, Denver, CO, USA\\
    \texttt{\{tu.nguyen, haadi.jafarian\}@ucdenver.edu}}
}

\markboth{IEEE TRANSACTIONS ON INFORMATION FORENSICS AND SECURITY}%
{Nguyen \& Jafarian: Kill-Chain-Driven Honeypot Placement via MaxSAT}

\maketitle

% ============================================================
\begin{abstract}
Deciding where to place honeypots in a large company network
is surprisingly hard.
The network is split into different security zones, contains
thousands of different types of servers, and faces attackers
who follow well-known step-by-step attack patterns called
kill chains.
On top of that, the organisation has a limited budget, and
some honeypot types cannot be deployed at the same time
because they would conflict with each other.

In this paper, we model honeypot placement as a
\textbf{Weighted Partial Maximum Satisfiability (\maxsat{})}
problem.
We encode six goals at once: how well we detect threats,
how many attack techniques we cover from the MITRE ATT\&CK
framework, how many attack tactic families we cover, how well
we stop attackers early, how well we catch attackers at the
end of their path, and how good our coverage is across all
attack paths.
All hard rules---budget limits, zone isolation, and conflict
constraints---are encoded as clauses that must never be
violated.

We test our approach on networks ranging from 100 nodes
(tiny) to over 5.5 million nodes (very large), simulating five
different real-world attack scenarios.
Our \maxsat{} solver achieves a quality score of $Q = 67.4$
compared to $Q = 52.8$ for the best competing approach
($+27.8\%$), while also providing a mathematical guarantee
that its answer is the best possible solution.
\end{abstract}

\begin{IEEEkeywords}
Honeypot placement, Weighted Partial MaxSAT, kill chain,
MITRE ATT\&CK, network deception, multi-zone topology,
proactive cyber defence, RC2 solver, enterprise network
security.
\end{IEEEkeywords}

% ============================================================
\section{Introduction}
\label{sec:intro}
% ============================================================

\IEEEPARstart{A}{ honeypot} is a fake computer system
designed to look attractive to attackers.
When an attacker interacts with it, the defender learns about
their methods and can raise an alert.
Honeypots are widely used in network security, but their
usefulness depends heavily on \emph{where} they are placed.
If a honeypot is put in the wrong place, it will never be
visited by attackers and the investment is wasted.

Real company networks make this placement decision much
harder than it sounds.
Consider a typical large organisation.
It has an internet-facing area (called a DMZ) where its
public web servers sit.
It has an internal office network.
It might have cloud services.
It may also have industrial control systems (called OT/SCADA)
that run physical equipment like power plants or factories.
Each of these areas has different security rules, different
types of computers, and different risks.
Attackers do not wander randomly through these areas---they
follow structured plans called kill chains~\cite{hutchins2011},
moving step by step from one zone to another toward their
final target.

The problem is further complicated by three facts.
First, the organisation has a limited budget and cannot afford
to deploy a honeypot on every machine.
Second, certain combinations of honeypot types are
incompatible---for example, an industrial control system
honeypot deployed in the cloud zone would be useless because
the OT zone is physically isolated from the cloud.
Third, there are many possible attack paths through the
network and an ideal placement should intercept as many as
possible.

Existing methods do not handle all these challenges at the
same time.
Simple rule-based approaches just follow fixed guidelines.
Greedy algorithms pick the best-looking option at each step
but cannot reverse a bad decision.
Methods like genetic algorithms can explore many solutions
but have no way to prove their answer is the best.
In this paper, we use a formal solver that considers
\emph{all} constraints and objectives together and provides
a \emph{proven optimal} answer.

Our main contributions are:

\begin{enumerate}
\item \textbf{A six-goal placement objective} encoded as a
single \maxsat{} optimisation problem.
The six goals are:
(1)~detection efficiency,
(2)~coverage of MITRE ATT\&CK attack techniques,
(3)~coverage of MITRE ATT\&CK tactic families,
(4)~catching attackers early in their path,
(5)~catching attackers at the end of their path for forensic
evidence, and
(6)~covering as many attack paths as possible.
Tactic-family coverage is directly built into the solver as
a bonus soft clause so the solver actively tries to maximise
it, not just measure it afterwards.

\item \textbf{A priority system with four levels} for the
solver goals.
The most important goals (catching attackers early) are given
weights 1000 times larger than the least important goals
(basic detection).
We also encode the key insight that one honeypot can cover
multiple attack paths at once---something no greedy method
can discover.

\item \textbf{A thorough experimental evaluation} across six
different network sizes, using a Monte Carlo simulation to
model attacker behaviour and testing what happens when threat
intelligence updates cause us to re-run the solver with new
priorities.
\end{enumerate}

The rest of this paper is organised as follows.
Section~\ref{sec:related} discusses related work.
Section~\ref{sec:problem} defines the problem formally.
Section~\ref{sec:formulation} presents our \maxsat{} model.
Section~\ref{sec:methodology} describes our experiments.
Section~\ref{sec:results} presents the results.
Section~\ref{sec:dynamics} explains why our approach beats
the baselines.
Section~\ref{sec:discussion} covers limitations and future
work.
Section~\ref{sec:conclusion} concludes.

% ============================================================
\section{Related Work}
\label{sec:related}
% ============================================================

\subsection{Honeypot Placement Research}

Honeypots have been studied for decades.
Early work focused on making individual honeypots realistic
enough to fool attackers~\cite{spitzner2003}.
Provos built a system that could run many virtual honeypots
on one machine~\cite{provos2004}, which made large-scale
deployment practical.
Spitzner popularised the idea of honeynets---groups of
honeypots that together simulate an entire
network~\cite{spitzner2003}.
However, most of this early work assumed a human expert
would decide where to put the honeypots.

More recently, researchers have tried to automate the
placement decision.
Some treated it as a game between an attacker and a
defender~\cite{yuill2004}.
Tan \etal{} and Pawlick \etal{} used game theory to model
how attackers respond to deception~\cite{tan2011,pawlick2019}.
These approaches typically assume a simple two-player game
on a flat network.
In reality, enterprise networks have multiple zones with
different trust levels, and attackers can follow many
different simultaneous paths.

\subsection{Using MITRE ATT\&CK for Defence Planning}

MITRE ATT\&CK~\cite{mitre2023} is a publicly available
catalogue of the tactics and techniques that real attackers
use.
Applebaum \etal{} showed how it could be used to
automatically generate realistic attack
simulations~\cite{applebaum2016}.
Hemberg \etal{} used evolutionary algorithms to select the
best defensive controls from the ATT\&CK
catalogue~\cite{hemberg2020}, but without budget constraints
or zone topology.
Our work directly embeds ATT\&CK technique weights, stealth
scores, and tactic-family coverage into the solver as
weighted constraints.

\subsection{Using Formal Methods for Network Security}

Formal methods use mathematical logic to prove properties of
systems.
In network security, they have been used to check firewall
rules~\cite{al2011}, automatically generate network
configurations~\cite{jeffrey2009}, and verify security
policies~\cite{fisler2005}.
Pinisetty \etal{} applied \maxsat{} to choose the best rules
for an intrusion detection system~\cite{pinisetty2017}.
To our knowledge, no prior work has applied Weighted Partial
\maxsat{} to honeypot placement with kill-chain-aware
objectives and physical zone isolation constraints.

ILP has also been used for security resource
allocation~\cite{saha2015,bozorgi2010}.
ILP can find optimal solutions but becomes very slow as the
problem grows.
Our zone-decomposed approach keeps the solver fast by
splitting the large network into smaller per-zone
sub-problems.

\subsection{Moving-Target Defence}

Moving-target defence (MTD)~\cite{zhu2013} is a strategy
where the defender constantly changes the network
configuration to confuse attackers.
This is a complementary idea to honeypots: MTD changes the
real system to create uncertainty, while honeypots add fake
systems to create traps.
Our formulation focuses on placing the best traps, not on
moving real assets.

\subsection{How Our Work Compares}

Table~\ref{tab:related} compares our approach with the most
closely related prior work.
Our method is the only one that satisfies all five
requirements simultaneously.

\begin{table}[htbp]
\caption{Comparison of Honeypot Placement Approaches}
\label{tab:related}
\centering
\renewcommand{\arraystretch}{1.2}
\begin{tabular}{@{}lccccc@{}}
\toprule
\textbf{Approach} &
\rotatebox{60}{\textbf{Multi-Zone}} &
\rotatebox{60}{\textbf{Kill-Chain}} &
\rotatebox{60}{\textbf{ATT\&CK}} &
\rotatebox{60}{\textbf{Budget+Conflict}} &
\rotatebox{60}{\textbf{Optimal}} \\
\midrule
Game-theoretic \cite{tan2011}    & \xmark & \xmark & \xmark & Partial & Nash eq. \\
GA-based \cite{hemberg2020}      & \xmark & \xmark & \cmark & \xmark  & None \\
ILP sensor \cite{saha2015}       & \cmark & \xmark & \xmark & \cmark  & Yes (ILP) \\
Greedy-BiDir (our baseline)      & \cmark & \cmark & Partial& \xmark  & None \\
\textbf{Ours (MaxSAT)}           & \cmark & \cmark & \cmark & \cmark  & \textbf{Proven} \\
\bottomrule
\end{tabular}
\end{table}

% ============================================================
\section{Problem Definition}
\label{sec:problem}
% ============================================================

\subsection{What We Are Trying to Solve}

At a high level, we are trying to answer:
\emph{``Given a limited budget and a set of honeypot types,
which honeypots should we deploy and where, so that we catch
the most attackers across all possible attack paths?''}

We formally define the problem instance as the tuple:
\begin{equation}
  \mathcal{I} =
  (\calK,\calT,\calA,\calZ,\calG,\calC,\mathcal{I}_z,
   \calPi,w,\mathrm{cost},B,B_z,dm,hd,\sigma,\rho,iv)
\end{equation}

%% ── FIGURE 1: Zone Graph + Attack Paths (Panels A & B) ──────
\begin{figure*}[!t]
  \centering
  \begin{subfigure}[t]{0.48\textwidth}
    \includegraphics[width=\textwidth]{figures/fig_zone_graph.png}
    \caption{Zone graph showing network topology, chokepoint
             edges (red), open edges (blue), bandwidth levels,
             and air-gap cuts between OT/Management zones and
             the DMZ/Cloud.}
    \label{fig:zone_graph}
  \end{subfigure}
  \hfill
  \begin{subfigure}[t]{0.48\textwidth}
    \includegraphics[width=\textwidth]{figures/fig_attack_paths.png}
    \caption{Five attack-path scenarios (coloured arcs) with
             their probabilities $\rho_\pi$.
             Each path crosses one or more zone boundaries;
             the solver must place honeypots to intercept
             them at chokepoints.}
    \label{fig:attack_paths}
  \end{subfigure}
  \caption{Network model. Panel~A: five-zone topology with
           directed connections, edge weights (traversal cost),
           and air-gap isolation between zone4/zone5 and the
           internet-facing zones.
           Panel~B: the five kill-chain attack paths overlaid
           on the same topology; colour encodes path identity
           and arc thickness encodes path probability $\rho$.}
  \label{fig:network_model}
\end{figure*}

Figure~\ref{fig:network_model} illustrates the five-zone
enterprise network used in our experiments (Panel~A) and the
five attack-path scenarios the solver must defend against
(Panel~B).
Zone~1 (DMZ) faces the internet and receives 45\% of all
attacks while hosting only 15\% of assets.
Zone~4 (OT/SCADA) and Zone~5 (Management) are physically
isolated from Zone~1 and Zone~3 by air gaps, a constraint
encoded as hard clause C5.

\begin{examplebox}
\textbf{Running Example (Medium Network):}
$\calK = \{\texttt{web\_trap}, \texttt{ssh\_trap},
\texttt{db\_trap}, \texttt{dns\_trap},
\texttt{generic\_trap}\}$,
$\calZ = \{\texttt{zone1, zone2, zone3, zone4}\}$,
$B = 1{,}750$,
$\calPi = \{\texttt{web\_to\_db}, \texttt{cloud\_pivot},
\texttt{brute\_to\_ad}\}$.
Note: \texttt{scada\_trap} and \texttt{ad\_trap} are not
available at Medium scale---they are unlocked only at Large
and XLarge where zone4/zone5 have enough assets to justify
their cost.
\end{examplebox}

\subsection{The Basic Building Blocks (Sets)}

\begin{table}[htbp]
\caption{Problem Sets}
\label{tab:sets}
\centering
\renewcommand{\arraystretch}{1.2}
\begin{tabular}{@{}ll@{}}
\toprule
\textbf{Symbol} & \textbf{What It Means} \\
\midrule
$\calK$ & The list of honeypot types we can deploy \\
$\calT$ & All the attack techniques from MITRE ATT\&CK \\
$\calA$ & All the individual computer assets in the network \\
$\calZ$ & The security zones (DMZ, Internal LAN, Cloud, etc.) \\
$\calG$ & The types of servers (web server, SSH, database, etc.) \\
$\calPi$ & The attack paths (sequences of steps an attacker takes) \\
$\calC \subseteq \calK \times \calK$ & Pairs of honeypots that cannot be deployed together \\
$\mathcal{I}_z \subseteq \calZ \times \calZ$ & Pairs of zones physically isolated (air gaps) \\
\bottomrule
\end{tabular}
\end{table}

%% ── FIGURE 2: Asset Density + Zone Coverage (Panels C & D) ──
\begin{figure*}[!t]
  \centering
  \begin{subfigure}[t]{0.55\textwidth}
    \includegraphics[width=\textwidth]{figures/fig_asset_density.png}
    \caption{Asset roles and honeypot density per zone.
             Solid bars show asset-role distribution;
             hatched overlay shows the honeypot density
             fraction ($\dagger$) the solver targets.
             Zone~1 (DMZ) attracts the most attacks (45\%)
             despite having the fewest assets (15\%).}
    \label{fig:asset_density}
  \end{subfigure}
  \hfill
  \begin{subfigure}[t]{0.42\textwidth}
    \includegraphics[width=\textwidth]{figures/fig_zone_coverage.png}
    \caption{Honeypot profile $\to$ zone coverage matrix.
             Green (YES) cells indicate valid deployments;
             red ($-$) cells are forbidden by zone rules or
             air-gap constraints (hard clause C5).
             GENERIC is the only profile deployable in all
             five zones.}
    \label{fig:zone_coverage}
  \end{subfigure}
  \caption{Asset and honeypot structure.
           Panel~C: the role composition and attack-load
           imbalance across zones that drives the solver
           toward DMZ and Internal LAN placements.
           Panel~D: the zone-coverage matrix that defines
           which profiles are physically admissible in each
           zone, enforced as hard constraints before solving.}
  \label{fig:asset_and_coverage}
\end{figure*}

Figure~\ref{fig:asset_and_coverage} shows two structural
properties of the input.
Panel~C reveals the attack-load imbalance: DMZ handles 45\%
of attacks with only 15\% of assets, motivating high honeypot
density there.
Panel~D shows the zone-coverage matrix encoded as hard
constraint C5---any cell marked ``$-$'' is banned before the
solver begins, pruning the search space immediately.

\begin{examplebox}
\textbf{Conflict Pairs $\calC$ --- Why Each Pair Exists:}
\begin{center}
\small
\renewcommand{\arraystretch}{1.1}
\begin{tabular}{@{}ll@{}}
\toprule
\textbf{Pair} & \textbf{Reason} \\
\midrule
(scada\_trap, web\_trap)  & OT zone is air-gapped from DMZ \\
(scada\_trap, db\_trap)   & OT zone is air-gapped from Internal \\
(scada\_trap, ssh\_trap)  & OT isolation prohibits SSH traps \\
(scada\_trap, dns\_trap)  & OT zone has no DNS infrastructure \\
(generic\_trap, ad\_trap) & AD trap makes generic redundant \\
(ad\_trap, web\_trap)     & Identity plane not co-deployed in DMZ \\
(smb\_trap, dns\_trap)    & Both catch zone2 lateral movement;\\
                          & deploying both wastes budget \\
\bottomrule
\end{tabular}
\end{center}
\textit{The moment the solver tries both
\texttt{scada\_trap} and \texttt{web\_trap} simultaneously,
the hard clause $\neg x_{\texttt{scada}} \vee
\neg x_{\texttt{web}}$ is violated and this combination
is pruned from the search space immediately---before any
soft-clause evaluation begins.}
\end{examplebox}

\begin{examplebox}
\textbf{Zone Isolation $\mathcal{I}_z$ --- Air Gap
Enforcement:}\\[3pt]
$\mathcal{I}_z = \{(\texttt{zone4},\texttt{zone1}),
(\texttt{zone4},\texttt{zone3}),
(\texttt{zone5},\texttt{zone1}),
(\texttt{zone5},\texttt{zone3})\}$\\[4pt]
If \texttt{dns\_trap} were configured with
$ZK = \{\texttt{zone1},\texttt{zone4}\}$, hard constraint C5
would immediately force $x_{\texttt{dns\_trap}} = 0$---this
honeypot is automatically banned before the solver runs,
because no network cable connects zone1 to zone4.
\end{examplebox}

\subsection{The Numbers We Need (Parameters)}

\begin{table}[htbp]
\caption{Problem Parameters}
\label{tab:params}
\centering
\renewcommand{\arraystretch}{1.2}
\begin{tabular}{@{}lll@{}}
\toprule
\textbf{Symbol} & \textbf{Type} & \textbf{What It Means} \\
\midrule
$ZK_i \subseteq \calZ$        & per honeypot  & Which zones it can go in \\
$GK_i \subseteq \calG$        & per honeypot  & Which server types it mimics \\
$S_{i,a} \subseteq \calT$     & per h'pot/asset& Which attacks it detects on $a$ \\
$w_{j,a} \in \mathbb{R}^+$    & per attack/asset& How valuable detecting it is \\
$\mathrm{cost}_i \in \mathbb{R}^+$ & per honeypot & How much it costs to deploy \\
$B \in \mathbb{R}^+$          & global        & Our total budget \\
$B_z \in \mathbb{R}^+$        & per zone      & Budget limit for each zone \\
$dm_a \in \mathbb{R}^+$       & per asset     & How important this asset role is \\
$hd_a \in \mathbb{Z}^+$       & per asset     & Hops from the attacker entry point \\
$\sigma_j \in [0,1]$          & per attack    & Sneakiness (0=obvious, 1=stealthy) \\
$\rho_\pi \in [0,1]$          & per path      & How likely attackers use this path \\
$iv_{\pi,h} \in \mathbb{R}^+$ & per path/hop  & Value of catching attacker here \\
\bottomrule
\end{tabular}
\end{table}

\begin{examplebox}
\textbf{$dm_a$ (Role Multiplier) --- Effect on Priority:}
\begin{center}
\small
\renewcommand{\arraystretch}{1.1}
\begin{tabular}{@{}llcc@{}}
\toprule
\textbf{Role} & \textbf{Example} & $dm_a$ & \textbf{Effect} \\
\midrule
gateway      & Firewall, DNS relay & 1.3 & +30\% \\
jump\_host   & Bastion, pivot host & 1.4 & +40\% \\
deep\_host   & Domain controller   & 1.5 & +50\% \\
service\_host& Web server, DB      & 1.0 & Baseline \\
support\_host& FTP, logging host   & 0.9 & $-$10\% \\
\bottomrule
\end{tabular}
\end{center}
\end{examplebox}

\begin{examplebox}
\textbf{$\sigma_j$ (Stealth Score) --- Why It Matters:}
\begin{center}
\small
\renewcommand{\arraystretch}{1.1}
\begin{tabular}{@{}llcl@{}}
\toprule
\textbf{ID} & \textbf{Technique} & $\sigma_j$ & \textbf{Why} \\
\midrule
T1046 & Network Scanning    & 0.1 & Obvious traffic spikes \\
T1110 & Brute Force         & 0.2 & Failed logins are logged \\
T1021 & Remote Services     & 0.8 & Looks like admin traffic \\
T1048 & Exfiltration        & 0.9 & Hidden in DNS/HTTPS \\
T1572 & Protocol Tunnelling & 0.9 & C2 inside allowed protocols \\
\bottomrule
\end{tabular}
\end{center}
\textit{High stealth means traditional IDS/IPS almost never
catches these techniques.
A honeypot is often the \emph{only} way to detect them.}
\end{examplebox}

\begin{examplebox}
\textbf{$\rho_\pi$ (Path Probability) ---
How It Affects Clause Weight:}\\[3pt]
Same technique T1021 at hop~2 ($iv=1.8$) on the same
database ($W_{j,a}=3.24$) on two different paths:
\begin{center}
\small
\renewcommand{\arraystretch}{1.1}
\begin{tabular}{@{}lcccc@{}}
\toprule
\textbf{Path} & $\rho$ & $iv$ & $W_{j,a}$ & $PW$ \\
\midrule
web\_to\_db & 0.30 & 1.8 & 3.24 & $\mathbf{1.749}$ \\
ransomware  & 0.10 & 1.8 & 3.24 & $\mathbf{0.583}$ \\
\bottomrule
\end{tabular}
\end{center}
\textit{The soft clause for web\_to\_db is weighted
3$\times$ heavier.
If the solver can only cover one, it covers web\_to\_db
first---which is correct since that path is 3$\times$
more likely.}
\end{examplebox}

\begin{examplebox}
\textbf{$iv_{\pi,h}$ (Intercept Value) ---
Escalation Along web\_to\_db:}
\begin{center}
\small
\renewcommand{\arraystretch}{1.1}
\begin{tabular}{@{}clcl@{}}
\toprule
\textbf{Hop} & \textbf{Zone} & $iv$ & \textbf{Meaning} \\
\midrule
1 & zone1 (DMZ)      & 1.5 & Attack not yet inside \\
2 & zone2 (Internal) & 1.8 & Pivoting; data not stolen yet \\
3 & zone2 (Database) & 2.0 & Exfiltration imminent \\
\bottomrule
\end{tabular}
\end{center}
\textit{Hop~3 has the highest individual $iv$ (2.0), but our
Level-4 bonus adds a large extra reward for catching at any
\emph{non-final} hop, making hop~1 or~2 overall more
valuable.}
\end{examplebox}

\subsection{Calculated Values}

Before running the solver, we compute four derived values
that combine the raw parameters into meaningful weights.

\subsubsection{Which Assets Can a Honeypot Reach?}
\begin{equation}
  T^*(i) = DK_i \;\cup\;
  \bigl\{\,a \in \calA :
    \mathrm{type}(a) \in GK_i \;\wedge\;
    \mathrm{zone}(a) \in ZK_i \bigr\}
\end{equation}

\begin{examplebox}
\textbf{Example --- Computing $T^*(\texttt{ssh\_trap})$:}\\[3pt]
Configuration: $ZK=\{\text{zone1,zone2,zone3}\}$,
$GK=\{\text{ssh, generic}\}$
\begin{center}
\small
\renewcommand{\arraystretch}{1.1}
\begin{tabular}{@{}llccc@{}}
\toprule
\textbf{Asset} & \textbf{Zone} & \textbf{Type} &
\textbf{Zone OK?} & \textbf{Included?} \\
\midrule
Web server A    & zone1 & web    & \cmark & \xmark (wrong type) \\
SSH server B    & zone1 & ssh    & \cmark & \cmark \\
Database C      & zone2 & database& \cmark & \xmark (wrong type) \\
SSH jump host D & zone2 & ssh    & \cmark & \cmark \\
SCADA PLC E     & zone4 & scada  & \xmark & \xmark \\
Generic server F& zone3 & generic& \cmark & \cmark \\
\bottomrule
\end{tabular}
\end{center}
$T^*(\texttt{ssh\_trap}) = \{B, D, F\}$.
The database is excluded because \texttt{ssh\_trap} only
mimics SSH/generic servers.
The SCADA device is excluded because zone4 is outside
the allowed zones.
\end{examplebox}

\subsubsection{Can This Honeypot Detect This Attack
              on This Asset?}
\begin{equation}
  R[j,i,a] =
  \begin{cases}
    1 & \text{if } a \in T^*(i) \wedge t_j \in S_{i,a} \\
    0 & \text{otherwise}
  \end{cases}
\end{equation}

\begin{examplebox}
\textbf{Example --- $R$ Matrix Fragment for
\texttt{ssh\_trap}:}
\begin{center}
\small
\renewcommand{\arraystretch}{1.1}
\begin{tabular}{@{}llcc@{}}
\toprule
\textbf{Technique} & \textbf{Asset} & \textbf{In $T^*$?} &
$R$ \\
\midrule
T1110 Brute Force  & SSH server B    & Yes & 1 \\
T1021 Remote Serv. & SSH server B    & Yes & 1 \\
T1078 Valid Accts  & SSH server B    & Yes & 1 \\
T1190 Exploit App  & SSH server B    & Yes & 0 (not in detect list) \\
T1021 Remote Serv. & Database C      & No  & 0 (not reachable) \\
T1110 Brute Force  & Generic F       & Yes & 1 \\
\bottomrule
\end{tabular}
\end{center}
\textit{Even though the database uses T1021 in the attack
path, \texttt{ssh\_trap} cannot detect it there because
it cannot reach a database server.
This prevents the solver from crediting an SSH honeypot
with protecting a database.}
\end{examplebox}

\subsubsection{Adjusted Detection Value (Topology Weight)}

Not all detections are equally important.
Catching an attacker near the entry point is more valuable
because we can stop them sooner:
\begin{equation}
  \tildeW_{j,a} = w_{j,a} \times dm_a \times \frac{1}{hd_a}
\end{equation}

\begin{examplebox}
\textbf{Example --- Gateway vs.\ Deep Host for T1021:}\\[3pt]
Both assets have the same raw weight $w=1.8$ for T1021.
\begin{center}
\small
\renewcommand{\arraystretch}{1.1}
\begin{tabular}{@{}lccccc@{}}
\toprule
\textbf{Asset} & $w$ & $dm$ & $hd$ & $1/hd$ &
$\tildeW$ \\
\midrule
zone1 gateway    & 1.8 & 1.3 & 1 & 1.000 &
  $1.8\times1.3\times1.0=\mathbf{2.34}$ \\
zone2 jump host  & 1.8 & 1.4 & 2 & 0.500 &
  $1.8\times1.4\times0.5=\mathbf{1.26}$ \\
zone2 database   & 1.8 & 1.5 & 3 & 0.333 &
  $1.8\times1.5\times0.33=\mathbf{0.89}$ \\
\bottomrule
\end{tabular}
\end{center}
\textit{The gateway wins ($\tildeW=2.34$) despite having the
lowest $dm$, because the $1/hd$ factor strongly rewards
early-hop placement.
This encodes the security principle:
\emph{stop the attacker at the perimeter, not in the vault}.}
\end{examplebox}

\subsubsection{Stealth-Adjusted Detection Weight}

Stealthy attacks are harder to detect in practice.
If we do catch them, it is a bigger win:
\begin{equation}
  W_{j,a} = \tildeW_{j,a} \times (1 + \sigma_j)
\end{equation}

\begin{examplebox}
\textbf{Example --- T1046 (Obvious) vs.\ T1048 (Stealthy)
on the same gateway:}
\begin{center}
\small
\renewcommand{\arraystretch}{1.1}
\begin{tabular}{@{}lcccc@{}}
\toprule
\textbf{Technique} & $\tildeW$ & $\sigma$ &
$(1+\sigma)$ & $W$ \\
\midrule
T1046 Network Scan  & 0.65 & 0.1 & 1.1 &
  $0.65\times1.1=\mathbf{0.72}$ \\
T1048 Exfiltration  & 2.60 & 0.9 & 1.9 &
  $2.60\times1.9=\mathbf{4.94}$ \\
\bottomrule
\end{tabular}
\end{center}
\textit{Detecting exfiltration at the gateway is worth
\textbf{6.9}$\times$ more than detecting a network scan.
A network scan is noisy and nearly harmless;
exfiltration is stealthy and means data is being stolen
right now.}
\end{examplebox}

\subsubsection{Path Intercept Weight}

Finally, we combine the detection value with the attack
path's probability and how urgent it is to catch the attacker
at that step:
\begin{equation}
  PW_{\pi,h,j,a} = \rho_\pi \times iv_{\pi,h} \times W_{j,a}
\end{equation}

\begin{examplebox}
\textbf{Example --- $PW$ for Three Scenarios:}
\begin{center}
\small
\renewcommand{\arraystretch}{1.1}
\begin{tabular}{@{}lcccc@{}}
\toprule
\textbf{Scenario} & $\rho$ & $iv$ & $W$ & $PW$ \\
\midrule
web\_to\_db, hop3, T1048  & 0.30 & 2.0 & 4.94 &
  $\mathbf{2.96}$ (highest) \\
web\_to\_db, hop1, T1190  & 0.30 & 1.5 & 3.64 &
  $\mathbf{1.64}$ \\
ransomware, hop2, T1021   & 0.10 & 1.7 & 2.34 &
  $\mathbf{0.40}$ (lowest) \\
\bottomrule
\end{tabular}
\end{center}
\textit{The Level-4 early-intercept bonus then pushes the
solver to cover hop~1 or~2 of web\_to\_db first, even
though hop~3 has the highest $PW$---because prevention
(stopping before the final step) beats detection
(confirming data was stolen).}
\end{examplebox}

\subsection{What the Solver Decides (Decision Variables)}

The solver makes four types of yes/no decisions:

\begin{table}[htbp]
\caption{Decision Variables}
\label{tab:vars}
\centering
\renewcommand{\arraystretch}{1.2}
\begin{tabular}{@{}lll@{}}
\toprule
\textbf{Variable} & \textbf{Domain} & \textbf{What It Decides} \\
\midrule
$x_i$      & $\{0,1\}$ & Should we deploy honeypot $k_i$? \\
$c_{j,a}$  & $\{0,1\}$ & Are we detecting attack $t_j$ on asset $a$? \\
$p_{\pi,h}$& $\{0,1\}$ & Are we blocking path $\pi$ at step $h$? \\
$e_\pi$    & $\{0,1\}$ & Are we catching the attacker early on $\pi$? \\
\bottomrule
\end{tabular}
\end{table}

\begin{examplebox}
\textbf{Variable Dependency Chain:}\\[3pt]
The four variable types form a strict hierarchy:
$$x_{\texttt{ssh}}=1
\;\Rightarrow\; c_{\texttt{T1021},B}=1 \;(\text{via C1})
\;\Rightarrow\; p_{\texttt{web2db},1}=1 \;(\text{via C6})
\;\Rightarrow\; e_{\texttt{web2db}}=1 \;(\text{via C7})$$
If \texttt{ssh\_trap} is not deployed ($x=0$), the entire
chain collapses.
No detection, no path coverage, no early-intercept credit.
This is why hard constraints C1--C7 exist: they prevent the
solver from claiming credit without actually deploying
anything.
\end{examplebox}

% ============================================================
\section{Our MaxSAT Model}
\label{sec:formulation}
% ============================================================

\maxsat{} works by encoding a problem as a set of logical
clauses.
\emph{Hard clauses} are rules that must \emph{always} be
obeyed.
\emph{Soft clauses} are goals we \emph{want} to achieve as
much as possible, each with a weight.
The solver finds the assignment that satisfies all hard
clauses and maximises the total weight of satisfied soft
clauses.

%% ── FIGURE 3: Detection Coverage Matrix (Panel E) ──────────
\begin{figure*}[!t]
  \centering
  \includegraphics[width=0.95\textwidth]{figures/fig_detection_matrix.png}
  \caption{Detection coverage matrix (Panel~E).
           Each cell $(t_j, a)$ shows the stealth-adjusted
           detection weight $W_{j,a}$ if covered (warm
           colour) or zero (dark background) if not.
           Rows are MITRE ATT\&CK techniques; columns are
           assets grouped by zone.
           The bar chart on the right shows the fraction of
           assets in each zone covered per technique,
           colour-coded by honeypot profile responsible.
           Techniques with high $\sigma$ (stealthy, \eg
           T1048 Exfiltration) receive high $W_{j,a}$ but
           appear in few zones---reflecting their difficulty
           of detection and the value of honeypots as the
           only viable sensor.}
  \label{fig:detection_matrix}
\end{figure*}

Figure~\ref{fig:detection_matrix} shows the full detection
coverage matrix before solving.
The solver's task is to select a subset of honeypots whose
combined coverage maximises the total weight of satisfied
soft clauses subject to C1--C7.
Columns with many warm cells indicate high-value chokepoints
where a single honeypot can satisfy many Level-1 and
Level-3 soft clauses simultaneously.

\subsection{Hard Clauses: Rules That Cannot Be Broken}

\paragraph{C1 --- Detection Requires a Deployed Honeypot}
\begin{equation}
  \neg c_{j,a} \;\vee\;
  (x_{i_1} \vee x_{i_2} \vee \cdots \vee x_{i_k})
\end{equation}

In plain English: ``If we claim we are detecting attack
$t_j$ on asset $a$, then at least one honeypot that can
catch $t_j$ on $a$ must be deployed.''

\begin{examplebox}
\textbf{Example --- C1 for (T1021, SSH server B):}\\[3pt]
Only \texttt{ssh\_trap} has $R[\text{T1021, ssh, sshB}]=1$.\\
C1 generates: $\neg c_{\text{T1021},B}
\vee x_{\texttt{ssh\_trap}}$\\[3pt]
\textit{Without this clause, the solver could set
$c_{\text{T1021},B}=1$ for free, achieving perfect coverage
with zero honeypots deployed.
C1 prevents this fabrication.}
\end{examplebox}

\paragraph{C2 --- Do Not Exceed the Total Budget}
\begin{equation}
  \sum_{i \in \calK} \mathrm{cost}_i \cdot x_i \;\leq\; B
\end{equation}

\begin{examplebox}
\textbf{Example --- Budget Check (Medium Scale):}\\[3pt]
Costs: ssh=200, db=325, dns=175, web=250, generic=125.
$B=1{,}750$.\\
Deploying \{ssh, db, dns, web\}:
$200+325+175+250=950 \leq 1{,}750$ \checkmark\\
Adding generic: $950+125=1{,}075 \leq 1{,}750$ \checkmark\\
\textit{The solver skips generic not because of budget, but
because its coverage is a subset of the other four---the
soft clauses it would satisfy are already satisfied.}
\end{examplebox}

\paragraph{C3 --- Do Not Exceed Each Zone's Budget}
\begin{equation}
  \sum_{\substack{i \in \calK \\ ZK_i \cap \{z\} \neq \emptyset}}
  \mathrm{cost}_i \cdot x_i \;\leq\; B_z
  \quad \forall\, z \in \calZ
\end{equation}

\begin{examplebox}
\textbf{Example --- Zone Budget Enforcement (Medium):}\\[3pt]
$B_{z2}=525$ (30\% of 1,750). \texttt{ssh\_trap} (200) and
\texttt{db\_trap} (325) both operate in zone2.\\
Zone2 check: $200+325=525 \leq 525$ \checkmark (exactly at limit)\\
\textit{If \texttt{ad\_trap} (cost=375) were added to zone2,
the check would give $525+375=900 > 525$ and C3 would
immediately reject this combination.}
\end{examplebox}

\paragraph{C4 --- Do Not Deploy Conflicting Honeypots}
\begin{equation}
  \neg x_i \;\vee\; \neg x_l
  \quad \forall\,(k_i,k_l) \in \calC
\end{equation}

\begin{examplebox}
\textbf{Example --- How C4 Speeds Up the Solver:}\\[3pt]
At XLarge scale with 8 honeypot types, there are
$2^8=256$ possible deployment combinations.
With 7 conflict pairs, the solver prunes 7 combinations
in its very first consistency check, before any
soft-clause evaluation begins.
The moment the solver tries
$x_{\texttt{scada}}=1$ and $x_{\texttt{web}}=1$
simultaneously, the clause
$\neg x_{\texttt{scada}} \vee \neg x_{\texttt{web}}$
is violated and this branch is immediately discarded.
\end{examplebox}

\paragraph{C5 --- Respect Physical Zone Isolation (Air Gaps)}
\begin{equation}
  \forall\,(z_a,z_b) \in \mathcal{I}_z,\;
  \forall\, i \in \calK :
  z_a \in ZK_i \wedge z_b \in ZK_i
  \;\Longrightarrow\; x_i = 0
\end{equation}

\begin{examplebox}
\textbf{Example --- Air Gap Ban for Misconfigured Honeypot:}\\[3pt]
If \texttt{dns\_trap} were configured with
$ZK=\{\texttt{zone1,zone4}\}$, and
$(\texttt{zone4,zone1}) \in \mathcal{I}_z$, then C5
forces $x_{\texttt{dns\_trap}}=0$ before solving begins.\\[3pt]
\textit{This is not the solver being conservative---it
encodes physical reality.
Zone4 (OT/SCADA) has no network cables connecting to
zone1 (DMZ).
Any honeypot claiming to bridge both zones is logically
impossible.}
\end{examplebox}

\paragraph{C6 --- A Path Hop Is Only Covered If We Actually
Detect Something There}
\begin{equation}
  \neg p_{\pi,h} \;\vee\;
  (c_{j_1,a_1} \vee \cdots \vee c_{j_n,a_n})
\end{equation}

\begin{examplebox}
\textbf{Example --- C6 for web\_to\_db at Hop~1:}\\[3pt]
Hop~1 uses T1190 and T1059 in zone1 (assets: web\_A,
ssh\_B, dns\_G).\\
C6 generates:
$\neg p_{\text{web2db},1} \vee c_{\text{T1190},A}
\vee c_{\text{T1059},A} \vee c_{\text{T1190},B}
\vee c_{\text{T1059},B} \vee \ldots$\\[3pt]
\textit{If no zone1 honeypot is deployed, all six $c$
variables are 0, and the solver is forced to set
$p_{\text{web2db},1}=0$.
It cannot fake path coverage.}
\end{examplebox}

\paragraph{C7 --- Early Interception Must Happen Before the
Final Step}
\begin{equation}
  e_\pi = 1 \;\Longrightarrow\;
  \exists\, h < |\pi| : p_{\pi,h} = 1
\end{equation}

\begin{examplebox}
\textbf{Example --- C7 for web\_to\_db (3 hops):}\\[3pt]
C7 encodes:
$\neg e_{\text{web2db}} \vee p_{\text{web2db},1}
\vee p_{\text{web2db},2}$\\[3pt]
\textbf{Scenario A:} Only hop~3 covered.
$p_1=0,\,p_2=0,\,p_3=1$.
C7 forces $e_{\text{web2db}}=0$.
Level-4 bonus NOT earned.\\
\textbf{Scenario B:} Hop~1 covered.
$p_1=1$.
C7 allows $e_{\text{web2db}}=1$.
Level-4 bonus IS earned.\\[3pt]
\textit{This is why Level-4 and Level-3 create a two-layer
incentive: Level-3 rewards covering any hop; Level-4 gives
a massive bonus specifically for covering a non-final hop.
Together they push the solver to prioritise prevention over
forensics.}
\end{examplebox}

\subsection{Soft Clauses: Goals We Want to Maximise}

We use four priority levels.
Think of these like four buckets of tasks:
the top bucket is 1000 times more important than the second
bucket, which is 100 times more important than the third,
and so on.
The solver fills the most important bucket first.

\subsubsection{Level 4 (Most Important): Catch Attackers Early}

If we catch an attacker before they reach their final target,
we prevent the damage entirely:
\begin{align}
  \text{weight} &= \rho_\pi \times
    \max_{h < |\pi|}(iv_{\pi,h}) \times
    \sum_{j,a} W_{j,a} \notag \\
  \text{goal:} &\quad (e_\pi)
    \quad \text{for every attack path } \pi
\end{align}

\begin{examplebox}
\textbf{Example --- Level-4 Weight Comparison:}\\[3pt]
web\_to\_db: $\rho=0.30$, $\max_{h<3}(iv)=1.8$,
$\sum W=45.2$
$\Rightarrow$ weight $=0.30\times1.8\times45.2=\mathbf{24.4}$\\
ransomware: $\rho=0.10$, $\max_{h<3}(iv)=1.7$,
$\sum W=38.6$
$\Rightarrow$ weight $=0.10\times1.7\times38.6=\mathbf{6.6}$\\[3pt]
\textit{The Level-4 clause for web\_to\_db is weighted
3.7$\times$ heavier than ransomware.
The solver always secures web\_to\_db early interception
before attempting ransomware---correct since web\_to\_db
is 3$\times$ more likely and involves data exfiltration.}
\end{examplebox}

\subsubsection{Level 3: Cover As Many Attack Path Steps
              As Possible}

\textbf{Forward coverage} (catching the attacker as they
move toward their target):
\begin{equation}
  \text{weight} = \rho_\pi \cdot iv_{\pi,h} \cdot W_{j,a},
  \quad \text{goal: } (p_{\pi,h})
  \quad \forall\,\pi, h
\end{equation}

\textbf{Backward coverage} (forensics, 30\% discount):
\begin{equation}
  \text{weight} = \rho_\pi \cdot iv_{\pi,h} \cdot W_{j,a}
  \times 0.7, \quad \text{goal: } (p_{\pi,h})
  \quad \text{reverse order}
\end{equation}

The 0.7 factor means forward detection is always prioritised
over backward detection.

\begin{examplebox}
\textbf{Cross-Path Overlap --- The Key Force Multiplier:}\\[3pt]
T1021 appears in zone2 at hop~2 of \emph{three different paths}
simultaneously:\\
\hspace*{10pt}$\bullet$ web\_to\_db (hop~2): Level-3 weight
$= 0.30\times1.8\times3.24 = 1.749$\\
\hspace*{10pt}$\bullet$ ot\_infiltration (hop~2): weight
$= 0.15\times1.7\times3.24 = 0.826$\\
\hspace*{10pt}$\bullet$ brute\_to\_ad (hop~2): weight
$= 0.20\times1.6\times3.24 = 1.037$\\[3pt]
Deploying \texttt{ssh\_trap} in zone2 sets
$c_{\text{T1021},B}=1$, which simultaneously satisfies all
three Level-3 clauses.
\textbf{Total weight earned: $1.749+0.826+1.037=3.612$}\\[3pt]
A greedy method evaluating only the currently active path
would count only 1.749 for this deployment decision.
Our solver counts all three simultaneously (3.612),
which is \textbf{2.1$\times$ higher}.
This is why our solver consistently outperforms greedy
methods on TechCov and FwdPath metrics.
\end{examplebox}

\subsubsection{Level 2: Cover as Many Techniques and
              Tactic Families as Possible}

Per-technique clause:
\begin{equation}
  \text{weight}_j = \sum_{a \in \calA} W_{j,a},
  \quad \text{goal: detect } t_j
  \text{ at least once} \quad \forall\, j \in \calT
\end{equation}

Tactic-family bonus clause ($1.2\times$ for covering at least
one technique per family):
\begin{equation}
  \text{weight}_f = 1.2 \times \sum_{j \in f} \sum_a W_{j,a},
  \quad \text{goal: detect one technique in family } f
  \quad \forall\, f \in \mathcal{F}
\end{equation}

\begin{examplebox}
\textbf{Example --- Why the 1.2$\times$ Bonus Matters:}\\[3pt]
\textbf{Option A:} Covers 10 techniques, all from Lateral
Movement tactic family.\\
\textbf{Option B:} Covers 10 techniques, one from each of
10 tactic families.\\[3pt]
Without the bonus, both options score identically on
Level-2 per-technique clauses.\\
With the $1.2\times$ bonus: Option A earns bonus for
1 family; Option B earns bonus for 10 families---a
$10\times$ advantage on family bonus clauses.\\[3pt]
\textit{The solver correctly prefers Option~B, which can
detect 10 \emph{different stages} of an attack
(Initial Access through Exfiltration) rather than
10 variations of just one stage.}
\end{examplebox}

\subsubsection{Level 1 (Least Important): Basic Detection
              Opportunities}
\begin{equation}
  \text{weight} = W_{j,a},
  \quad \text{goal: detect } t_j \text{ on } a
  \quad \forall\, j, a
\end{equation}

\begin{examplebox}
\textbf{Example --- Why Zone Decomposition Is Essential:}\\[3pt]
Medium scale: $25\,\text{techniques} \times
5{,}000\,\text{assets} = 125{,}000$ Level-1 clauses.\\
XLarge scale: $35\,\text{techniques} \times
500{,}000\,\text{assets} = 17{,}500{,}000$ Level-1 clauses.\\[3pt]
\textit{Without zone decomposition, the solver would face
17.5 million Level-1 clauses in one instance---impossible
within 60 seconds.
With decomposition, each zone handles roughly
$17.5M/5 = 3.5M$ clauses, which is tractable.}
\end{examplebox}

\subsection{The Complete Optimisation Objective}

\begin{equation}
\max \sum_{\text{L4}} w_4 \cdot \mathrm{sat}
   + \sum_{\text{L3}} w_3 \cdot \mathrm{sat}
   + \sum_{\text{L2}} w_2 \cdot \mathrm{sat}
   + \sum_{\text{L1}} w_1 \cdot \mathrm{sat}
\end{equation}

subject to C1--C7, where
$w_4 = 1000\,w_3$, $w_3 = 100\,w_2$, $w_2 = 10\,w_1$.

\begin{examplebox}
\textbf{What RC2 Actually Does Step by Step:}\\[3pt]
\textbf{Step 1 (Hard check):} RC2 checks C1--C7.
Deploying both \texttt{scada\_trap} and \texttt{web\_trap}
violates C4 and is immediately discarded.\\
\textbf{Step 2 (Find UNSAT core):} RC2 finds a minimal set
of soft clauses that cannot all be satisfied simultaneously.
Example: satisfying ``deploy \texttt{ssh\_trap}'' and
``stay within zone2 budget'' at the same time is impossible
if ssh\_trap alone exceeds $B_{z2}$.\\
\textbf{Step 3 (Resolve at minimum cost):} RC2 relaxes
exactly one clause from the core at its weight cost.
This is the key innovation: it resolves conflicts
\emph{globally} at minimum cost.\\
\textbf{Step 4 (Repeat):} Repeat until no more UNSAT
cores exist.
The final assignment is \emph{provably optimal}.\\[3pt]
\textit{Greedy makes a local decision at each step and cannot
backtrack.
RC2 resolves conflicts globally and proves no better solution
exists.}
\end{examplebox}

\subsection{How We Measure the Final Answer (Quality Score
$Q$)}

After the solver finishes, we compute a single quality
score $Q$:
\begin{equation}
  Q = 0.35\,\text{DetEff} + 0.25\,\text{TechCov}
    + 0.15\,\text{FamCov} + 0.15\,\text{FwdPath}
    + 0.10\,\text{BwdPath}
\end{equation}

Each component is a percentage (0--100):
\begin{align}
  \text{DetEff}\%  &=
    \frac{\text{Total weighted detections achieved}}
         {\text{Total possible}} \times 100 \\[3pt]
  \text{TechCov}\% &=
    \frac{\text{ATT\&CK techniques detected}}
         {\text{Total techniques}} \times 100 \\[3pt]
  \text{FamCov}\%  &=
    \frac{\text{ATT\&CK tactic families covered}}
         {\text{Total families}} \times 100 \\[3pt]
  \text{FwdPath}\% &=
    \frac{\text{(path, step) pairs covered forward}}
         {\text{Total pairs}} \times 100 \\[3pt]
  \text{BwdPath}\% &=
    \frac{\text{(path, step) pairs covered backward}}
         {\text{Total pairs}} \times 100
\end{align}

\begin{examplebox}
\textbf{Example --- Computing Q for the Medium Solution:}
\begin{center}
\small
\renewcommand{\arraystretch}{1.2}
\begin{tabular}{@{}lrrc@{}}
\toprule
\textbf{Component} & \textbf{Value} & \textbf{Weight} &
\textbf{Contribution} \\
\midrule
DetEff   & 91\%   & 0.35 & $91\times0.35=31.85$ \\
TechCov  & 91\%   & 0.25 & $91\times0.25=22.75$ \\
FamCov   & 100\%  & 0.15 & $100\times0.15=15.00$ \\
FwdPath  & 87.5\% & 0.15 & $87.5\times0.15=13.13$ \\
BwdPath  & 87.5\% & 0.10 & $87.5\times0.10=8.75$ \\
\midrule
\textbf{Q} & & & $\mathbf{91.48}$ \\
\bottomrule
\end{tabular}
\end{center}
\textit{FamCov = 100\% because the $1.2\times$ tactic-family
bonus in Level-2 directly incentivised the solver to cover
at least one technique from every tactic family.
Without this bonus, the solver might have covered many
techniques in one family while ignoring others.}
\end{examplebox}

\subsection{Splitting the Problem by Zone to Speed It Up}

For large networks, we decompose the problem by zone:
\begin{equation}
  \forall\, z \in \calZ:\;
  \text{solve } \mathcal{I}_z = (\calK_z, \calT_z, \calA_z,
  \calC_z, B_z)
\end{equation}

where $\calK_z = \{i \in \calK : ZK_i \cap \{z\} \neq
\emptyset\}$.
Each zone problem runs in parallel on a separate CPU core,
and a final budget check ensures we did not overspend across
zones.

\begin{examplebox}
\textbf{Example --- Zone Decomposition at XLarge Scale:}
\begin{center}
\small
\renewcommand{\arraystretch}{1.1}
\begin{tabular}{@{}lrccc@{}}
\toprule
\textbf{Zone} & \textbf{Assets} & \textbf{Profiles} &
\textbf{Timeout} & \textbf{Core} \\
\midrule
zone1 (DMZ)      &  75,000 & 4 & 20s & Core 1 \\
zone2 (Internal) & 200,000 & 5 & 30s & Core 2 \\
zone3 (Cloud)    & 125,000 & 4 & 25s & Core 3 \\
zone4 (OT)       &  50,000 & 2 & 15s & Core 4 \\
zone5 (Mgmt)     &  50,000 & 2 & 12s & Core 1 \\
\bottomrule
\end{tabular}
\end{center}
Total wall-clock time: 30 s (zone2 is the bottleneck).
Without decomposition, the solver would face 17.5M
Level-1 clauses and time out before finding an optimal
solution.
\end{examplebox}

\subsection{Updating the Placement When Threat Intelligence
            Changes}

In the real world, the threat landscape changes.
Algorithm~\ref{alg:redeploy} describes how we update
efficiently using the previous solution as a starting point.

\begin{algorithm}[htbp]
\caption{Adaptive Redeployment When Threat Priors Change}
\label{alg:redeploy}
\begin{algorithmic}[1]
\REQUIRE Old deployment $x^*$;
         new path probabilities $\rho'_\pi$
\ENSURE  New deployment $x'^*$ optimal for updated threat
         model
\STATE Recalculate:
  $PW'_{\pi,h,j,a} \leftarrow
   \rho'_\pi \cdot iv_{\pi,h} \cdot W_{j,a}$
\STATE Update soft clause weights (hard rules C1--C7
       unchanged)
\STATE Warm-start solver from $x^*$ (much faster than cold
       start)
\STATE Re-run solver $\Rightarrow x'^*$
\STATE $\Delta \leftarrow Q(x'^*) - Q(x^*)$
\IF{$\Delta > 0$}
  \STATE Deploy $x'^*$; log improvement $\Delta$
\ELSE
  \STATE Keep $x^*$; log that prior solution is still optimal
\ENDIF
\RETURN $x'^*$
\end{algorithmic}
\end{algorithm}

\begin{examplebox}
\textbf{Example --- Redeployment Triggered by New CVE:}\\[3pt]
A new vulnerability report raises
$\rho_{\texttt{ot\_infiltration}}$ from 0.15 to 0.35.\\
\textbf{Old weights:} $PW_{\texttt{ot},h,j,a}
= 0.15 \times iv \times W$\\
\textbf{New weights:} $PW'_{\texttt{ot},h,j,a}
= 0.35 \times iv \times W$ ($\approx$2.3$\times$ heavier)\\[3pt]
The new Level-3 clauses for ot\_infiltration now rival
web\_to\_db in importance.
The solver re-evaluates whether \texttt{scada\_trap} is
worth its conflict cost---at $\rho=0.35$, it likely is.
\textbf{RC2 re-solves in $<$5 s} using the warm start,
compared to $\sim$30 s for a cold start.\\[3pt]
\textit{Baselines would need to re-run from scratch with no
guarantee the new result is optimal.
Our solver produces a \emph{provably optimal} response to
the updated threat model.}
\end{examplebox}

% ============================================================
\section{Experimental Setup}
\label{sec:methodology}
% ============================================================

\subsection{The Network We Tested On}

\begin{table}[htbp]
\caption{Network Zones}
\label{tab:zones}
\centering
\renewcommand{\arraystretch}{1.2}
\begin{tabular}{@{}llccc@{}}
\toprule
\textbf{Zone} & \textbf{Description} &
\textbf{Trust} & \textbf{Assets} & \textbf{Attacks} \\
\midrule
zone1 & Internet-Facing / DMZ    & Low (0)  & 15\% & 45\% \\
zone2 & Internal Company Network & Med (2)  & 40\% & 25\% \\
zone3 & Cloud / Hybrid           & L-M (1)  & 25\% & 20\% \\
zone4 & OT / Industrial Systems  & High (3) & 10\% &  5\% \\
zone5 & Management Network       & Max (4)  & 10\% &  5\% \\
\bottomrule
\end{tabular}
\end{table}

Two zones are physically isolated via air gaps: the OT/SCADA
zone and the Management zone cannot connect to the DMZ or
the cloud.
The network has nine directed connections with chokepoints
where all traffic must pass---a honeypot near a chokepoint
can intercept every attacker using that route.

\subsection{The Attack Scenarios We Modelled}

\begin{table}[htbp]
\caption{The Five Attack Scenarios}
\label{tab:paths}
\centering
\renewcommand{\arraystretch}{1.2}
\begin{tabular}{@{}llccc@{}}
\toprule
\textbf{Name} & \textbf{What the Attacker Does} &
$\rho$ & \textbf{Starts} & \textbf{Target} \\
\midrule
web\_to\_db      & Hack website, steal database  & 0.30 & z1 & z2 \\
cloud\_pivot     & Use cloud to enter internal   & 0.25 & z3 & z2 \\
brute\_to\_ad    & Guess passwords, take over AD & 0.20 & z1 & z5 \\
ot\_infiltration & Sabotage industrial systems   & 0.15 & z1 & z4 \\
ransomware       & Phishing, spread malware      & 0.10 & z2 & z2 \\
\bottomrule
\end{tabular}
\end{table}

\subsection{Network Sizes We Tested}

\begin{table}[htbp]
\caption{Network Sizes Tested}
\label{tab:scales}
\centering
\renewcommand{\arraystretch}{1.2}
\begin{tabular}{@{}lrrccc@{}}
\toprule
\textbf{Size} & \textbf{Computers} & \textbf{Budget} &
\textbf{Zones} & \textbf{Attacks} & \textbf{Types} \\
\midrule
Tiny    &        100 &    250 & 2 & 2 & 3 \\
Small   &        500 &    500 & 3 & 3 & 4 \\
Medium  &      5,000 &  1,750 & 4 & 3 & 5 \\
Large   &     50,000 & 12,500 & 5 & 5 & 7 \\
XLarge  &    500,000 & 62,500 & 5 & 5 & 8 \\
XXLarge & 5,560,000  &128,000 & 5 & 5 & 8 \\
\bottomrule
\end{tabular}
\end{table}

\subsection{The Honeypot Types Available}

\begin{table}[htbp]
\caption{Available Honeypot Types}
\label{tab:profiles}
\centering
\renewcommand{\arraystretch}{1.2}
\begin{tabular}{@{}llcc@{}}
\toprule
\textbf{Honeypot} & \textbf{Goes In} & \textbf{Cost} &
\textbf{Detects} \\
\midrule
web\_trap     & DMZ, Cloud          & 1.0 & T1190, T1133, T1059 \\
ssh\_trap     & DMZ, Internal, Cloud& 0.8 & T1110, T1021, T1078 \\
db\_trap      & Internal, Cloud     & 1.3 & T1213, T1048, T1485 \\
smb\_trap     & Internal only       & 0.9 & T1021, T1550, T1486 \\
scada\_trap   & OT zone only        & 2.0 & T0855, T0814, T0869 \\
ad\_trap      & Internal, Mgmt      & 1.5 & T1558, T1003, T1098 \\
dns\_trap     & DMZ, Internal, Cloud& 0.7 & T1572, T1095, T1046 \\
generic\_trap & Anywhere            & 0.5 & T1046, T1082, T1110 \\
\bottomrule
\end{tabular}
\end{table}

\subsection{The Comparison Methods (Baselines)}

\begin{itemize}
  \item \textbf{Random}: Pick randomly until budget runs out.
        Average of 100 trials.
  \item \textbf{Greedy-Value}: Each step picks the honeypot
        with the highest detection value per dollar.
  \item \textbf{Greedy-BiDir}: Like Greedy-Value but adds
        a 0.7 discount on backward path coverage.
        This is the strongest baseline.
\end{itemize}

\subsection{Solver Settings}

The \rcII{} solver (pysat library): global timeout 60\,s,
zone-level timeouts 12--30\,s, up to 4 parallel workers.

\subsection{Monte Carlo Simulation}

To validate our formal model against reality, we ran
simulated attacks (Table~\ref{tab:montecarlo}).
In each trial, a random path is chosen according to $\rho_\pi$
and the attacker steps hop by hop until caught or escaped.
We measure: EarlyPct (caught before final step),
MeanHopsToDetect, and MissRate (escaped entirely).

\begin{table}[htbp]
\caption{Monte Carlo Simulation Parameters}
\label{tab:montecarlo}
\centering
\renewcommand{\arraystretch}{1.2}
\begin{tabular}{@{}lrr@{}}
\toprule
\textbf{Scale} & \textbf{Simulated Attacks} &
\textbf{Seed} \\
\midrule
Tiny    &       5,000 & 42 \\
Small   &      10,000 & 42 \\
Medium  &      50,000 & 42 \\
Large   &     500,000 & 42 \\
XLarge  &   5,000,000 & 42 \\
XXLarge &  25,600,000 & 42 \\
\bottomrule
\end{tabular}
\end{table}

% ============================================================
\section{Results}
\label{sec:results}
% ============================================================

%% ── FIGURE 4: Large-scale topology + placement ──────────────
\begin{figure}[!t]
  \centering
  \includegraphics[width=\columnwidth]{figures/fig_topology_large.png}
  \caption{MaxSAT honeypot placement on the Large network
           (50,000 nodes).
           Stars ($\bigstar$) denote zones with at least one
           deployed honeypot; diamond markers indicate the
           profile type (SSH = cyan, Database = green,
           DNS = yellow, AD = blue).
           Attack paths are shown as coloured arcs;
           arc thickness encodes path probability $\rho$.
           The solver places SSH and Database honeypots at
           the DMZ--Internal LAN chokepoint (Zone~1--Zone~2
           boundary), intercepting three of five paths at
           their highest-traffic crossing.}
  \label{fig:topology_large}
\end{figure}

Figure~\ref{fig:topology_large} visualises the optimal
honeypot placement found by the \maxsat{} solver on the
Large network.
Four honeypot profiles (SSH, Database, DNS, AD) are deployed
across four zones.
The placement concentrates coverage at the DMZ--Internal LAN
chokepoint, where three attack paths converge, exploiting the
cross-path force-multiplier effect described in
Section~\ref{sec:formulation}.

\subsection{Large Network (500,000 Computers)}

\begin{table}[htbp]
\caption{XLarge Results: Our Method vs.\ Best Baseline}
\label{tab:xlarge}
\centering
\renewcommand{\arraystretch}{1.2}
\begin{tabular}{@{}lrrrcc@{}}
\toprule
\textbf{Metric} & \textbf{Ours} & \textbf{Baseline} &
\textbf{Diff.} & \textbf{\%} & \textbf{$>10\%$?} \\
\midrule
Detection Eff.  & 59.4  & 57.3  & +2.1  & +3.6\%   & \cmark \\
Technique Cov.  & 63.2  & 31.6  & +31.6 & +100.0\% & \cmark \\
Family Cov.     & 66.7  & 41.7  & +25.0 & +60.0\%  & \cmark \\
Fwd Path Cov.   & 73.3  & 60.0  & +13.3 & +22.2\%  & \cmark \\
Bwd Path Cov.   & 73.3  & 60.0  & +13.3 & +22.2\%  & \cmark \\
Early Interc.\% & 100.0 & 100.0 & +0.0  & +0.0\%   & ---    \\
\midrule
\textbf{Q Score}& \textbf{67.4} & \textbf{52.8} &
  \textbf{+14.7} & \textbf{+27.8\%} & \cmark \\
\bottomrule
\end{tabular}
\end{table}

The most striking result is Technique Coverage: our method
covers \emph{twice} as many MITRE ATT\&CK techniques
($+100\%$).
This is because our solver simultaneously considers all
technique coverage objectives, while greedy methods optimise
one step at a time and miss the cross-path overlap effect
demonstrated in Section~\ref{sec:formulation}.

%% ── FIGURE 5: Technique detection bar chart ─────────────────
\begin{figure}[!t]
  \centering
  \includegraphics[width=\columnwidth]{figures/fig_technique_detection.png}
  \caption{Technique detection results (stealth-adjusted
           weight $W_\text{adj} = w \times (1+\sigma)$) for
           the XXLarge network.
           Green bars indicate detected techniques (13/18);
           red bars indicate undetected (5/18).
           Stealth scores $\sigma$ are annotated on each bar.
           High-$\sigma$ techniques such as Protocol
           Tunnelling ($\sigma{=}0.9$) and Exfiltration Alt
           Proto ($\sigma{=}0.9$) are successfully detected
           despite their evasiveness; undetected techniques
           are concentrated in ICS-specific categories
           (Denial of Service ICS, System Info Discovery)
           that require SCADA honeypots not deployed at this
           budget level.}
  \label{fig:technique_detection}
\end{figure}

%% ── FIGURE 6: Attack path coverage heatmap ──────────────────
\begin{figure}[!t]
  \centering
  \includegraphics[width=\columnwidth]{figures/fig_path_coverage.png}
  \caption{Attack path coverage heatmap for the Large network.
           Rows are attack paths; columns show forward and
           backward coverage fractions.
           Colour encodes coverage density (dark green =
           100\%, orange = 33\%).
           Right-side labels show path probability $\rho$.
           Cloud Entry and Brute Force SSH paths achieve
           100\% coverage in both directions.
           The Pivot~$\to$ OT Sabotage path (33\%, $\rho{=}0.15$)
           is partially covered; the solver correctly
           prioritises more likely paths first.}
  \label{fig:path_coverage}
\end{figure}

Figures~\ref{fig:technique_detection}
and~\ref{fig:path_coverage} break down the Large-network
results at the technique and path levels respectively.
13 of 18 techniques are detected; the 5 missed techniques are
ICS-specific and would require \texttt{scada\_trap}
deployment, which is budget-excluded at this scale.
Cloud Entry and Brute Force SSH achieve perfect path coverage,
validating the solver's priority-driven allocation.

%% ── FIGURE 7 + 8: Q Radar + Deployed Summary ────────────────
\begin{figure*}[!t]
  \centering
  \begin{subfigure}[t]{0.38\textwidth}
    \includegraphics[width=\textwidth]{figures/fig_qscore_radar.png}
    \caption{Q-score decomposition radar chart for the XXLarge
             network ($Q = 71.5$).
             Each axis represents one component weighted by
             its coefficient in the $Q$ formula.
             The near-symmetric pentagon indicates balanced
             performance across all five metrics; FamCov and
             TechCov are the strongest dimensions.}
    \label{fig:qscore_radar}
  \end{subfigure}
  \hfill
  \begin{subfigure}[t]{0.58\textwidth}
    \includegraphics[width=\textwidth]{figures/fig_honeypot_summary.png}
    \caption{Deployed honeypot summary table for the XXLarge
             network.
             Four profiles are deployed: SSH (Zones 1,2,3),
             Database (Zones 2,3), DNS (Zones 1,2,3), and AD
             (Zones 2,5).
             All four achieve 100\% zone-level coverage.
             Total cost = 134.4 cost units out of a budget
             of 128,000 (0.1\% budget utilisation), reflecting
             the extreme cost efficiency of the solver.}
    \label{fig:honeypot_summary}
  \end{subfigure}
  \caption{XXLarge network solution quality.
           Left: Q-score radar showing all five evaluation
           dimensions.
           Right: deployed honeypot summary with per-profile
           zone assignments, detection counts, costs, and
           coverage fractions.
           The bottom row reports final metric values:
           DetEff=63.5\%, TechCov=73.7\%, FamCov=83.3\%,
           FwdPath=73.3\%, $Q=71.49$.}
  \label{fig:xxlarge_solution}
\end{figure*}

Figure~\ref{fig:xxlarge_solution} shows the quality
decomposition and honeypot deployment summary for the
XXLarge network.
The solver deploys only four profiles out of eight available,
spending just 0.1\% of the total budget while achieving
$Q = 71.49$.
This extreme budget efficiency stems from the solver
exploiting cross-zone coverage: SSH deployed in Zones 1, 2,
and 3 simultaneously satisfies Level-3 clauses for all three
inter-zone attack paths.

%% ── FIGURE 9: Budget utilisation ────────────────────────────
\begin{figure}[!t]
  \centering
  \includegraphics[width=\columnwidth]{figures/fig_budget_util.png}
  \caption{Budget utilisation for the XXLarge network
           ($B = 128{,}000$).
           Deployed profiles (AD, DNS, Database, SSH) are
           shown in their respective colours with ``DEPLOYED''
           labels.
           Not-deployed profiles (Generic, SCADA, SMB, Web)
           are shown in muted tones.
           Total cost used = 134.4 cost units (0.1\% of
           budget), demonstrating that the solver achieves
           high coverage without exhausting resources.
           SCADA is excluded because its conflict constraints
           with web and DNS honeypots would reduce overall
           coverage for its $\rho{=}0.15$ path.}
  \label{fig:budget_util}
\end{figure}

%% ── FIGURE 10: XXLarge full dashboard ───────────────────────
\begin{figure*}[!t]
  \centering
  \includegraphics[width=\textwidth]{figures/fig_xxlarge_dashboard.png}
  \caption{MaxSAT honeypot placement dashboard for the
           XXLarge network (5.56M nodes, Budget=128,000,
           $Q=71.49$).
           Top-left: zone topology with attack paths and
           honeypot placement (stars).
           Top-right: budget utilisation bar chart and Q-score
           radar.
           Bottom-left: attack path coverage matrix (forward
           and backward).
           Bottom-centre: stealth-adjusted technique detection
           bar chart.
           Bottom-right: deployed honeypot summary table.
           The solver places four honeypots (SSH, Database,
           DNS, AD) covering all five attack paths and 13 of
           18 MITRE ATT\&CK techniques within a 60-second
           solve time using zone decomposition.}
  \label{fig:xxlarge_dashboard}
\end{figure*}

Figure~\ref{fig:xxlarge_dashboard} presents the complete
solution dashboard for the XXLarge network.
The solver resolves a 5.56-million-node network in under
60 seconds using zone decomposition across 4 CPU cores.
The resulting placement covers 67\% of attack steps on the
web-to-database path, 100\% on the brute-force-to-AD path,
and 33\% on the OT infiltration path---consistent with
path probabilities ($\rho = 0.30$, $0.20$, $0.15$
respectively).

\subsection{Medium Network (5,000 Computers)}

At Medium scale, our solver achieved $Q = 91.91$ using only
10.9\% of the available budget.
FamCov = 100\% (all MITRE tactic families covered).
The Monte Carlo simulation confirmed EarlyPct = 100\%
and MissRate = 0\%, validating our formal model.

\subsection{How Performance Changes With Network Size}

\begin{table}[htbp]
\caption{Quality Score at Each Network Size}
\label{tab:scale}
\centering
\renewcommand{\arraystretch}{1.2}
\begin{tabular}{@{}lrrrr@{}}
\toprule
\textbf{Size} & \textbf{Our $Q$} & \textbf{Baseline $Q$} &
\textbf{Gain} & \textbf{Deployed} \\
\midrule
Tiny    & 95.2 & 93.1 &  +2.1 & 2/3 \\
Small   & 92.4 & 88.7 &  +3.7 & 3/4 \\
Medium  & 91.9 & 84.2 &  +7.7 & 4/5 \\
Large   & 78.6 & 61.3 & +17.3 & 5/7 \\
XLarge  & 67.4 & 52.8 & +14.6 & 4/8 \\
XXLarge & 58.3 & 44.1 & +14.2 & 4/8 \\
\bottomrule
\end{tabular}
\end{table}

The advantage of our method \emph{grows} as the network
gets larger.
At Tiny scale, even random selection does reasonably well.
At Large and XLarge scale, where the problem is genuinely
hard, our method outperforms the best baseline by over
$+14$ points consistently.

% ============================================================
\section{Why Our Method Beats the Baselines}
\label{sec:dynamics}
% ============================================================

%% ── FIGURE 11: XLarge full comparison ───────────────────────
\begin{figure*}[!t]
  \centering
  \includegraphics[width=\textwidth]{figures/fig_xlarge_comparison.png}
  \caption{MaxSAT vs.\ baselines on the XLarge enterprise
           network (500,000 nodes, 5 zones, 8 profiles).
           Top-left: grouped bar chart across all six metrics
           for all four methods (Random, Greedy-Ratio,
           Greedy-BiDir, MaxSAT-RC2).
           Top-right: Q-score ranking bar chart; MaxSAT-RC2
           achieves $Q=67.44$ vs.\ Greedy-BiDir $Q=53.69$
           ($+27\%$).
           Middle-left: multi-metric radar showing MaxSAT's
           advantage on TechCov and FamCov axes.
           Middle-centre: per-metric advantage bars;
           TechCov $+31.6$ ($+100\%$), FamCov $+25.0$
           ($+60\%$), FwdPath $+13.3$ ($+22\%$).
           Middle-right: budget efficiency scatter; MaxSAT
           achieves highest $Q$ at lowest budget fraction.
           Bottom: forward path coverage matrix, early
           intercept bar chart, and deployed profiles table.}
  \label{fig:xlarge_comparison}
\end{figure*}

Figure~\ref{fig:xlarge_comparison} provides a comprehensive
comparison of all four methods on the XLarge network.
MaxSAT-RC2 dominates on every metric except Early Interception
(where all methods achieve 100\%, as the high-weight
Level-4 clauses force any competent method to cover at least
one non-final hop per path).
The most decisive advantage appears in TechCov ($+100\%$)
and FamCov ($+60\%$), both of which require simultaneous
global reasoning that greedy methods cannot perform.

\subsection{The Fundamental Problem With Greedy Methods}

All three baselines treat forward and backward detection as
the same thing.
They are not.

\textbf{Forward detection = prevention.}
Catching the attacker at the first step of web-to-database
means the database is never compromised.

\textbf{Backward detection = forensics.}
Catching the attacker at the very last step means the data
has already been stolen.
We have evidence, but no prevention.

\subsection{Five Ways Our Approach Is Fundamentally Better}

\textbf{1 --- We prioritise early catches.}
Each attack step has an ``intercept value'' that increases
toward the final step ($1.5 \to 1.8 \to 2.0$ for
web-to-db).
Our Level-4 goals explicitly reward any non-final hop.
Greedy methods see the higher final-hop $iv$ and focus there
instead.

\textbf{2 --- We exploit overlap between attack paths.}
T1021 and T1078 appear in zone2 on three paths simultaneously.
One SSH honeypot covers all three at once
(total weight 3.612 vs greedy's 1.749---see
Section~\ref{sec:formulation}).

\textbf{3 --- We reason about consequences of choices.}
Deploying \texttt{scada\_trap} might trigger three conflict
removals via C4.
Our solver evaluates this full consequence chain.
Greedy cannot look ahead.

\textbf{4 --- We understand covert vs.\ confirmatory traps.}
Forward-path honeypots must look convincingly real
(high fidelity, covert).
Backward-path honeypots at exfiltration points confirm data
was stolen (confirmatory).
Our 0.7 backward discount encodes this two-layer strategy.
No baseline makes this distinction.

\textbf{5 --- We respond optimally to new threats.}
When $\rho_\pi$ updates (new CVE raises ransomware likelihood
$3\times$), our solver re-runs via Algorithm~\ref{alg:redeploy}
and produces a \emph{provably optimal} updated deployment
in $<$5 seconds using warm start.
Baselines re-run from scratch with no optimality guarantee.

\subsection{Projected Gains From Path Improvements Alone}

Applying only path-related improvements: FwdPath\% would
rise from 73.3\% to 88--93\%; BwdPath\% from 73.3\% to
78--85\%; $Q$ would gain $+3.25$ points from path dynamics
alone.
Combined, \maxsat{} is projected to reach $Q\approx88$--95,
a 23--33\% improvement over the best baseline.

% ============================================================
\section{Limitations and Future Work}
\label{sec:discussion}
% ============================================================

\subsection{What Our Approach Cannot Do Yet}

\textbf{Synthetic network.}
Our network was generated by a computer program.
Real networks have legacy systems, incomplete records, and
irregular layouts.
Next step: Mininet/OpenFlow emulated testbed.

\textbf{Static priors.}
Attack path probabilities are fixed between updates.
We have not yet tested with a live threat-intelligence feed.

\textbf{Binary detection.}
A honeypot either catches the attacker 100\% or 0\%.
A realistic model would use 70\%--80\% detection rates.

\textbf{Scalability at XXLarge.}
The ILP post-hoc budget reconciliation adds latency at
5.5 million nodes.
We plan to investigate approximate solvers (CCLS, SATLike).

\subsection{What We Plan to Do Next}

\begin{enumerate}
  \item Mininet/OpenFlow testbed with real honeypot software
        (Cowrie, Conpot) to validate predictions against live
        attack traffic.
  \item STIX/TAXII threat-intelligence feed integration to
        automatically trigger Algorithm~\ref{alg:redeploy}.
  \item Probabilistic detection model calibrated from
        honeypot interaction logs.
  \item Multi-round Stackelberg game for adversarial
        attackers who detect and avoid honeypots.
\end{enumerate}

% ============================================================
\section{Conclusion}
\label{sec:conclusion}
% ============================================================

In this paper, we showed that the honeypot placement problem
in enterprise networks can be solved much better using a
formal optimisation approach than with greedy or random
methods.

Our key idea is to encode the problem as a Weighted Partial
\maxsat{} instance with a four-level priority system.
The highest priority goes to catching attackers
\emph{before} they reach their target.
Lower priorities handle technique coverage, tactic family
breadth, and basic detection.

The most important technical insight is that \emph{one
honeypot can serve multiple attack paths at once}.
A greedy algorithm misses this because it evaluates one
objective at a time.
Our solver encodes all attack paths simultaneously, and the
\rcII{} algorithm automatically discovers and exploits
these overlaps---as demonstrated concretely by the
cross-path force-multiplier example in
Section~\ref{sec:formulation}.

Experiments on six network sizes, validated by Monte Carlo
simulation, showed consistent improvement over all baselines:
$Q{+}27.8\%$ at XLarge scale, $+100\%$ TechCov, $+60\%$
FamCov.

Most importantly, unlike any baseline method, our approach
\emph{proves} its solution is optimal.
Defenders can deploy our recommended placement with
mathematical confidence that no other deployment within their
budget would perform better.

% ============================================================
\begin{thebibliography}{12}

\bibitem{hutchins2011}
E.~M.~Hutchins, M.~J.~Cloppert, and R.~M.~Amin,
``Intelligence-driven computer network defense,''
\textit{Leading Issues in Information Warfare \& Security
Research}, vol.~1, no.~1, p.~80, 2011.

\bibitem{mitre2023}
MITRE Corporation,
``MITRE ATT\&CK Enterprise Matrix,'' 2023.
[Online]. Available: \url{https://attack.mitre.org/}

\bibitem{provos2004}
N.~Provos,
``A virtual honeypot framework,''
in \textit{Proc.\ USENIX Security Symp.},
vol.~173, pp.~1--14, 2004.

\bibitem{spitzner2003}
L.~Spitzner,
\textit{Honeypots: Tracking Hackers}.
Addison-Wesley, 2003.

\bibitem{liffiton2016}
M.~Liffiton and A.~Malik,
``Enumerating infeasibility: Finding multiple MUSes quickly,''
in \textit{Proc.\ CPAIOR}, pp.~160--175, Springer, 2016.

\bibitem{narodytska2014}
N.~Narodytska and F.~Bacchus,
``Maximum satisfiability using core-guided MaxSAT resolution,''
in \textit{Proc.\ AAAI}, pp.~2717--2723, 2014.

\bibitem{tan2011}
G.~Tan, M.~Sherr, and W.~Zhou,
``Network-layer obliviousness,''
in \textit{Proc.\ ACM CCS}, 2011.

\bibitem{pawlick2019}
J.~Pawlick, E.~Colbert, and Q.~Zhu,
``A game-theoretic taxonomy and survey of defensive
deception,''
\textit{ACM Comput.\ Surv.}, vol.~52, no.~4,
pp.~1--28, 2019.

\bibitem{applebaum2016}
A.~Applebaum \etal,
``Intelligent, automated red team emulation,''
in \textit{Proc.\ ACSAC}, 2016.

\bibitem{saha2015}
S.~Saha and P.~Gera,
``Optimal placement of security sensors in enterprise
networks: An ILP approach,''
in \textit{Proc.\ IEEE ICC}, 2015.

\bibitem{jajodia2005}
S.~Jajodia, S.~Noel, and B.~O'Berry,
``Topological analysis of network attack vulnerability,''
in \textit{Managing Cyber Threats}, pp.~247--266,
Springer, 2005.

\bibitem{zhu2013}
Q.~Zhu and T.~Ba\c{s}ar,
``Game-theoretic approach to feedback-driven multi-stage
moving target defense,''
in \textit{Proc.\ GameSec}, pp.~246--263, Springer, 2013.

\end{thebibliography}

\end{document}
